This section defines how requirements and rules in this document shall be interpreted and prioritized.

\subsection{Requirement Characteristics}

Project requirements should be:
\begin{itemize}
    \item clear and unambiguous;
    \item necessary for the intended purpose;
    \item feasible within the available resources;
    \item verifiable through review, testing, analysis, or inspection;
    \item consistent with other applicable requirements;
    \item traceable to an identified source or stakeholder need.
\end{itemize}

When a requirement is unclear, the ambiguity shall be resolved before the affected work is finalized. The resolution should be recorded when it may affect implementation, testing, review, or future maintenance.

\subsection{Rule Priority}

When two requirements or rules conflict, they shall normally be applied in the following order of priority:
\begin{enumerate}
    \item applicable laws, regulations, and safety requirements;
    \item contractual, customer, or externally imposed requirements;
    \item approved project-specific requirements;
    \item requirements in this document;
    \item organizational or team conventions;
    \item general recommendations and personal preferences.
\end{enumerate}

A more specific rule takes precedence over a more general rule at the same priority level. A language-specific rule takes precedence over a general code rule when the two rules address the same subject.

Conflicts shall not be resolved silently. The conflict, its impact, and the chosen resolution should be documented when the decision may affect the project or future contributors.

\subsection{Mandatory and Advisory Rules}

Rules using the word ``\must'' are mandatory. A mandatory rule shall be followed unless an approved exception has been granted.

Rules using the word ``\should'' are advisory but represent the expected default practice. A deviation should have a reasonable technical
justification, especially when it affects maintainability, portability, security, or verification.

Rules using the word ``\may'' identify an allowed option and do not impose a requirement.

\subsection{Exceptions}

An exception is required when a mandatory rule cannot be followed. Exceptions shall be documented before release of the affected deliverable and should include:
\begin{itemize}
    \item the rule being waived;
    \item the affected files, components, or deliverables;
    \item the reason for the exception;
    \item the risks introduced by the deviation;
    \item any compensating controls;
    \item the approving person or authority;
    \item an expiration date or review condition, where appropriate.
\end{itemize}

Exceptions should be limited in scope and duration. Temporary exceptions shall be reviewed periodically and removed when they are no longer required.

\subsection{Compliance}

Compliance may be demonstrated through code review, automated checks, static analysis, testing, inspection, documentation review, or other appropriate verification methods.

A violation discovered during review shall be corrected, formally accepted as an exception, or recorded as a known issue with an assigned owner and planned resolution.